\section{Singular Vector Alignment Regularization for Model Merging}
\label{sec:gargiulo_regularization}

\subsection{Method}

We propose a regularization technique during fine-tuning that encourages the preservation of the principal subspace structure of pretrained weights. The key insight is that successful model merging benefits from maintaining alignment between the singular vector bases of fine-tuned models and their shared pretrained initialization.

For each weight matrix $W \in \mathbb{R}^{m \times n}$ in the encoder, we compute its truncated singular value decomposition (SVD):
\begin{equation}
W = U \Sigma V^\top \approx U_k \Sigma_k V_k^\top
\end{equation}
where $U_k \in \mathbb{R}^{m \times k}$ and $V_k \in \mathbb{R}^{n \times k}$ contain the top-$k$ left and right singular vectors, respectively.

The Gargiulo regularization penalty encourages the current model's singular vectors to remain aligned with those of the pretrained model. For the left singular vectors, the penalty is defined as:
\begin{equation}
\mathcal{L}_{\text{Gargiulo}} = \sum_{l} \left\| I_k - U_{\text{pre}}^{(l)\top} U_{\text{cur}}^{(l)} \right\|_F^2
\end{equation}
where $U_{\text{pre}}^{(l)}$ and $U_{\text{cur}}^{(l)}$ are the top-$k$ left singular vectors of layer $l$ for the pretrained and current model, respectively, and $I_k$ is the $k \times k$ identity matrix.

The total training objective combines the task-specific cross-entropy loss with the regularization term:
\begin{equation}
\mathcal{L}_{\text{total}} = \mathcal{L}_{\text{CE}} + \lambda \cdot \mathcal{L}_{\text{Gargiulo}}
\end{equation}

\subsection{Implementation Details}

We apply the Gargiulo regularization to all 2D weight matrices in the ViT-B-16 image encoder. The SVD of the current model weights is recomputed every 50 training steps to balance computational cost with regularization accuracy. The hyperparameters used in our experiments are:
\begin{itemize}
    \item Number of singular vectors: $k = 10$
    \item Regularization strength: $\lambda = 10^{-4}$
    \item SVD update interval: 50 steps
    \item Singular vector type: left singular vectors ($U$)
\end{itemize}

\subsection{Results}

We evaluate the effect of Gargiulo regularization on multi-task model merging using the ViT-B-16 architecture across two benchmark settings: 8 tasks (N8) and 20 tasks (N20). We compare five merging methods: arithmetic mean (Task Arithmetic), DARE, Task-Specific Vectors (TSV), TIES, and simple weight averaging.

\subsubsection{8-Task Benchmark (N8)}

Table~\ref{tab:gargiulo_n8} presents the normalized accuracy results for the 8-task benchmark.

\begin{table}[h]
\centering
\caption{Normalized accuracy on the 8-task benchmark (N8). $\Delta$ indicates absolute percentage point change from baseline.}
\label{tab:gargiulo_n8}
\begin{tabular}{lccc}
\toprule
\textbf{Merger} & \textbf{Baseline} & \textbf{Gargiulo} & \textbf{$\Delta$} \\
\midrule
Arithmetic & 0.7834 & 0.7930 & \textbf{+0.96\%} \\
DARE       & 0.7873 & 0.7895 & +0.22\% \\
TSV        & 0.9219 & 0.9219 & +0.00\% \\
TIES       & 0.7627 & 0.7661 & +0.34\% \\
Weight Avg & 0.7878 & 0.7890 & +0.12\% \\
\midrule
\textbf{Average} & \textbf{0.8086} & \textbf{0.8119} & \textbf{+0.33\%} \\
\bottomrule
\end{tabular}
\end{table}

On the 8-task benchmark, Gargiulo regularization provides consistent improvements across all merging methods except TSV (which already achieves near-optimal performance). The arithmetic merging method benefits most significantly (+0.96\%), followed by TIES (+0.34\%) and DARE (+0.22\%). The average improvement across all methods is +0.33 percentage points.

\subsubsection{20-Task Benchmark (N20)}

Table~\ref{tab:gargiulo_n20} presents the results for the more challenging 20-task benchmark.

\begin{table}[h]
\centering
\caption{Normalized accuracy on the 20-task benchmark (N20). $\Delta$ indicates absolute percentage point change from baseline.}
\label{tab:gargiulo_n20}
\begin{tabular}{lccc}
\toprule
\textbf{Merger} & \textbf{Baseline} & \textbf{Gargiulo} & \textbf{$\Delta$} \\
\midrule
Arithmetic & 0.4916 & 0.5033 & \textbf{+1.17\%} \\
DARE       & 0.7531 & 0.7516 & $-$0.15\% \\
TSV        & 0.8551 & 0.8551 & +0.00\% \\
TIES       & 0.7598 & 0.7586 & $-$0.13\% \\
Weight Avg & 0.7527 & 0.7529 & +0.02\% \\
\midrule
\textbf{Average} & \textbf{0.7225} & \textbf{0.7243} & \textbf{+0.18\%} \\
\bottomrule
\end{tabular}
\end{table}

On the 20-task benchmark, the regularization shows more nuanced effects. The arithmetic merging method continues to benefit substantially (+1.17\%), showing the largest improvement across both benchmarks. However, DARE and TIES exhibit slight performance degradation ($-$0.15\% and $-$0.13\%, respectively), while TSV remains unchanged and weight averaging shows marginal improvement (+0.02\%). The overall average improvement is +0.18 percentage points.

\subsection{Discussion}

The experimental results reveal several insights about singular vector alignment regularization:

\begin{enumerate}
    \item \textbf{Arithmetic merging benefits most}: The task arithmetic method consistently shows the largest improvements from Gargiulo regularization across both benchmarks (+0.96\% on N8, +1.17\% on N20). This suggests that simple averaging-based merging is particularly sensitive to the geometric alignment of weight spaces.

    \item \textbf{Diminishing returns with task complexity}: The overall benefit of regularization decreases as the number of tasks increases (from +0.33\% on N8 to +0.18\% on N20). This may indicate that maintaining singular vector alignment becomes more challenging when merging many diverse task-specific models.

    \item \textbf{Method-specific interactions}: Advanced merging methods like DARE and TIES, which incorporate sparsification or sign-based filtering, show reduced or negative effects from the regularization on the 20-task benchmark. This suggests that these methods may have internal mechanisms that interact non-trivially with the regularization.

    \item \textbf{TSV saturation}: The Task-Specific Vectors method remains unaffected by the regularization, likely because it already achieves near-optimal performance through its explicit handling of task-specific components.
\end{enumerate}

These findings suggest that Gargiulo regularization is most effective for simpler merging strategies and smaller-scale multi-task scenarios, while more sophisticated merging algorithms may require alternative or complementary regularization approaches.
